%!TEX root = ../QuantumNoiseAndDecoherence.tex
% ──────────────────────────────────────────────────────────────
%  Lecture 8 — Lindblad Equations II: Radiative Damping and
%              Cavity Emission
% ──────────────────────────────────────────────────────────────

% Continues Section 5 (Lindblad Equations) from Lecture 7.

\subsection{Recap: the radiative-damping model}%
\label{sec:radiative-damping-recap}

We continue the derivation begun in Lecture~7.  The setting is a
two-level atom (system~$S$) coupled to a bosonic field (environment~$E$)
via dipole interaction.  In the interaction picture, after the
\textbf{rotating-wave approximation} (RWA), the interaction Hamiltonian
is
\begin{equation}\label{eq:V-int-RWA}
  \hat{V}_{\mathrm{int}}(t)
  = \underbrace{\sum_k g_k\,\sigma_+\, \hat{b}_k\, e^{i(\omega_a - \omega_k)t}}_{V_+(t)}
  + \underbrace{\sum_k g_k^*\,\sigma_-\, \hat{b}_k^\adj\, e^{-i(\omega_a - \omega_k)t}}_{V_-(t)}\,,
\end{equation}
where $\sigma_\pm$ are the atomic raising/lowering operators,
$\hat{b}_k$, $\hat{b}_k^\adj$ are the field mode operators with
frequencies~$\omega_k$, and $g_k$ are the coupling constants.  The
initial state of the environment is taken to be the vacuum:
$\rho_E(0) = \ketbra{0}{0}$.

\begin{intuition}[Physical picture of radiative damping]
  With the field in the vacuum, no photons impinge on the atom.
  An excited atom can only lose energy: the term~$V_-(t)$ creates
  a photon in mode~$k$ while flipping the atom downward.  The created
  photon flies away and (by the Markov assumption) never returns.
  Hence an initial excited state $\ket{1}$ decays irreversibly to the
  ground state~$\ket{0}$.
\end{intuition}

% ──────────────────────────────────────────────────────────────
\subsection{Deriving the master equation}%
\label{sec:master-eq-derivation}

\subsubsection{The Born-approximation equation}

In the Born approximation the equation of motion for the system
density operator is
\begin{equation}\label{eq:born-eq}
  \frac{\dd}{\dd t}\,\rho_S(t)
  = -\int_0^t \dd t_1\;
    \tr_E\!\bigl[\hat{V}_{\mathrm{int}}(t),\,
    [\hat{V}_{\mathrm{int}}(t_1),\,\rho_S(t_1)\tp\ketbra{0}{0}]\bigr]\,.
\end{equation}
Expanding the double commutator with $V_{\mathrm{int}} = V_+ + V_-$
produces $4 \times 4 = 16$ terms.

\subsubsection{Vacuum simplification}

The key identity that eliminates most terms is:
\begin{equation}\label{eq:vacuum-identity}
  \tr_E\!\bigl((\One_S \tp A_E)\;\rho_S \tp \ketbra{0}{0}
    \;(\One_S \tp B_E)\bigr)
  = \bra{0}B_E\, A_E\ket{0}\;\rho_S\,,
\end{equation}
together with $\hat{b}_k\ket{0} = 0$ for all modes~$k$.  Since
$V_+(t)$ contains annihilation operators acting on the field, any term
where $V_+$ stands adjacent to $\ket{0}$ vanishes.  A non-zero
contribution requires equally many $V_+$ and $V_-$ factors, reducing
the 16~terms to just~4.

\begin{remark}
  The surviving terms pair creation and annihilation operators so
  that the field expectation value
  $\bra{0}\hat{b}_k\,\hat{b}_l^\adj\ket{0} = \delta_{kl}$ produces
  a Kronecker delta, collapsing double sums over modes to single sums.
\end{remark}

\subsubsection{Continuum limit and the bath correlation function}

Replacing the discrete mode sum by an integral introduces the
\textbf{density of states}~$\rho(\omega)$:
\begin{equation}\label{eq:continuum-limit}
  \sum_k \lvert g_k\rvert^2\, f(\omega_k)
  \;\longrightarrow\;
  \int_0^\infty \dd\omega\;\rho(\omega)\,\lvert g(\omega)\rvert^2\,f(\omega)\,.
\end{equation}
The bath correlation function that appears in~\eqref{eq:born-eq} is
\begin{equation}\label{eq:gamma-tau}
  \gamma(\tau)
  = \int_0^\infty \dd\omega\;
    \rho(\omega)\,\lvert g(\omega)\rvert^2\,
    e^{i(\omega_a - \omega)\tau}\,.
\end{equation}

\begin{keyresult}[Markov approximation]
  If the product $\rho(\omega)\lvert g(\omega)\rvert^2$ is
  \emph{slowly varying} around $\omega \approx \omega_a$, then its
  Fourier transform $\gamma(\tau)$ is sharply peaked near $\tau = 0$,
  acting effectively as a delta function under the time integral
  in~\eqref{eq:born-eq}.  This converts the integro-differential
  (non-Markovian) equation into a memoryless (Markovian) equation of
  motion.
\end{keyresult}

\subsubsection{The Lindblad equation for radiative damping}

Decomposing $\int_0^\infty \gamma(\tau)\,\dd\tau$ into real and
imaginary parts,
\begin{equation}\label{eq:gamma-decompose}
  \int_0^\infty \gamma(\tau)\,\dd\tau
  = \frac{\gamma}{2} + i\,\delta\omega_a\,,
\end{equation}
and combining all terms, the master equation becomes
\begin{equation}\label{eq:lindblad-rad-damping}
  \eqbox{%
  \frac{\dd\rho_S}{\dd t}
  = -i\bigl[\tfrac{\omega_a + \delta\omega_a}{2}\,\sz,\;\rho_S\bigr]
    + \gamma\,\mathcal{D}[\sigma_-]\rho_S}\,,
\end{equation}
where the \textbf{dissipator} (or \textbf{Lindblad superoperator}) is
defined by
\begin{equation}\label{eq:dissipator}
  \eqbox{%
  \mathcal{D}[L]\rho
  \coloneqq L\,\rho\, L^\adj
    - \tfrac{1}{2}\,L^\adj L\,\rho
    - \tfrac{1}{2}\,\rho\, L^\adj L}\,.
\end{equation}

\begin{definition}[Lindblad equation]\label{def:lindblad}
  An equation of the form
  $\dot{\rho} = -i[H,\rho] + \sum_j \gamma_j\,\mathcal{D}[L_j]\rho$,
  with $\gamma_j \geq 0$ and $H = H^\adj$, is called a
  \textbf{Lindblad (or GKSL) master equation}.  The operators~$L_j$
  are called \textbf{Lindblad operators} (or \textbf{jump operators}).
\end{definition}

\begin{remark}
  The frequency shift $\delta\omega_a$ is the \textbf{Lamb shift}:
  the coupling to the vacuum field renormalises the atomic transition
  frequency.  For most purposes it simply redefines $\omega_a$ and
  can be absorbed by moving to a suitable rotating frame.
\end{remark}

\begin{keyresult}[Complete positivity]
  Any Lindblad equation generates a completely positive, trace-preserving
  (CPTP) semigroup $\rho(t) = e^{\mathcal{L}t}\rho(0)$.
  Conversely, for finite-dimensional systems, every continuous
  Markovian CPTP evolution is generated by a Lindblad equation.  This
  is the content of the \emph{Gorini--Kossakowski--Sudarshan--Lindblad
  theorem}.
\end{keyresult}

% ──────────────────────────────────────────────────────────────
\subsection{Solving the radiative-damping equation}%
\label{sec:solve-rad-damping}

Setting $\delta\omega_a = 0$ (absorbing the Lamb shift into the
rotating frame), we solve the Lindblad
equation~\eqref{eq:lindblad-rad-damping} with
$L = \sigma_- = \ketbra{0}{1}$.  Parameterise the qubit density
operator as
\begin{equation}\label{eq:rho-param}
  \rho_S(t)
  = \begin{pmatrix} 1 - p(t) & x(t) \\ x^*(t) & p(t) \end{pmatrix}\,,
\end{equation}
where $p(t)$ is the excited-state population and $x(t)$ encodes the
coherence.  Substituting into the Lindblad equation yields two
decoupled equations:
\begin{equation}\label{eq:ode-rad-damping}
  \dot{p} = -\gamma\, p\,,
  \qquad
  \dot{x} = -\tfrac{\gamma}{2}\,x\,.
\end{equation}

\begin{keyresult}[Solution: radiative damping of a two-level atom]
  The qubit density operator at time~$t$ is
  \begin{equation}\label{eq:rho-rad-solution}
    \eqbox{%
    \rho_S(t)
    = \begin{pmatrix}
        1 - p(0)\,e^{-\gamma t} & x(0)\,e^{-\gamma t/2} \\
        x^*(0)\,e^{-\gamma t/2} & p(0)\,e^{-\gamma t}
      \end{pmatrix}}\,.
  \end{equation}
  \begin{itemize}
    \item The excited-state population decays at rate~$\gamma$.
    \item The coherences decay at rate~$\gamma/2$ (half the population
      decay rate).
    \item The \textbf{steady state} is the ground state:
      $\rho_{\mathrm{ss}} = \lim_{t\to\infty}\rho_S(t) = \ketbra{0}{0}$.
  \end{itemize}
\end{keyresult}

\begin{intuition}[Dynamics on the Bloch sphere]
  On the Bloch sphere, every initial state spirals inward toward the
  north pole~$\ket{0}$.  The excited state~$\ket{1}$ (south pole)
  moves straight up along the $z$-axis; superposition states follow
  curved trajectories as the coherence decays faster than the
  population.  After a characteristic time $\sim 1/\gamma$ the atom
  is, with high probability, in its ground state.
\end{intuition}

% ──────────────────────────────────────────────────────────────
\subsection{Cavity emission}%
\label{sec:cavity-emission}

As a second example, we consider a single cavity mode (a harmonic
oscillator with annihilation operator~$a$ and frequency~$\omega_c$)
that leaks through an imperfect mirror into external radiation modes.

\begin{intuition}[The leaky cavity]
  A photon inside the cavity bounces between the mirrors.  Each
  reflection at the imperfect mirror has a small probability of
  transmission; the photon eventually escapes and is lost to the
  external field.  This is the bosonic analogue of radiative damping.
\end{intuition}

The total Hamiltonian is quadratic (field--field coupling):
\begin{equation}\label{eq:cavity-hamiltonian}
  H = \omega_c\, a^\adj a
    + \sum_k \omega_k\, \hat{b}_k^\adj \hat{b}_k
    + \sum_k \bigl(g_k\, a^\adj \hat{b}_k + g_k^*\, a\, \hat{b}_k^\adj\bigr)\,.
\end{equation}
The derivation follows the same steps as for radiative damping with
the substitution $\sigma_+ \to a^\adj$, $\sigma_- \to a$.

\subsubsection{Vacuum environment}

If the external field is in the vacuum, we obtain
\begin{equation}\label{eq:lindblad-cavity-vac}
  \frac{\dd\rho_S}{\dd t}
  = -i[\omega_c\, a^\adj a,\,\rho_S]
    + \kappa\,\mathcal{D}[a]\,\rho_S\,,
\end{equation}
where $\kappa > 0$ is the cavity decay rate.

\subsubsection{Thermal environment}

If the external field is in a thermal state at frequency~$\omega_c$
with mean photon number $\bar{n} = (e^{\beta\omega_c} - 1)^{-1}$,
an additional term accounts for incoherent photon injection:
\begin{equation}\label{eq:lindblad-cavity-thermal}
  \eqbox{%
  \frac{\dd\rho_S}{\dd t}
  = -i[\omega_c\, a^\adj a,\,\rho_S]
    + \kappa(\bar{n}+1)\,\mathcal{D}[a]\,\rho_S
    + \kappa\bar{n}\,\mathcal{D}[a^\adj]\,\rho_S}\,.
\end{equation}

% ──────────────────────────────────────────────────────────────
\subsection{Solving the cavity equation via the characteristic
function}\label{sec:cavity-char-fn}

For the vacuum-environment case~\eqref{eq:lindblad-cavity-vac}
(working in the frame rotating at~$\omega_c$), we translate the master
equation into an equation for the characteristic function
$\chi(\psi,t) = \tr\bigl(\rho_S(t)\, D(\psi)\bigr)$.  Using the
identities
\begin{equation}\label{eq:D-deriv-identities}
  a^\adj\, D(\psi)
  = \Bigl(\frac{\partial}{\partial\psi} + \frac{\psi^*}{2}\Bigr)D(\psi)\,,
  \qquad
  a\, D(\psi)
  = \Bigl(-\frac{\partial}{\partial\psi^*} + \frac{\psi}{2}\Bigr)D(\psi)\,,
\end{equation}
the master equation becomes a first-order PDE:
\begin{equation}\label{eq:chi-pde}
  \frac{\partial\chi}{\partial t}
  = -\frac{\kappa}{2}\Bigl(
      \psi\,\frac{\partial\chi}{\partial\psi}
    + \psi^*\frac{\partial\chi}{\partial\psi^*}
    + \lvert\psi\rvert^2\,\chi
    \Bigr)\,.
\end{equation}
This is a Fokker--Planck-type equation whose solution preserves the
Gaussian form.  Making the ansatz
$\chi(\psi,t) = \exp\!\bigl(-\tfrac{1}{2}K(t)\lvert\psi\rvert^2\bigr)$
for an initial Gaussian state with parameter~$K(0)$, one finds
\begin{equation}\label{eq:chi-solution}
  K(t)
  = 1 + \bigl(K(0) - 1\bigr)\,e^{-\kappa t}\,.
\end{equation}

\begin{keyresult}[Steady state of a leaky cavity]
  As $t \to \infty$, $K(t) \to 1$ and
  $\chi(\psi) \to e^{-\lvert\psi\rvert^2/2}$, which is the
  characteristic function of the vacuum state.  Every initial state of
  the cavity decays to the vacuum at rate~$\kappa$.
\end{keyresult}

\begin{exercise}\label{ex:lindblad-derivation}
  Complete the derivation of the remaining three terms in the radiative
  damping master equation~\eqref{eq:born-eq}.  Use the vacuum
  identities and verify that the result
  matches~\eqref{eq:lindblad-rad-damping}.
\end{exercise}

\begin{exercise}\label{ex:cavity-lindblad}
  Starting from the cavity Hamiltonian~\eqref{eq:cavity-hamiltonian}
  with a thermal environment, carry out the Born--Markov derivation
  to obtain~\eqref{eq:lindblad-cavity-thermal}.  Identify where the
  additional $\mathcal{D}[a^\adj]$ term arises.
\end{exercise}

\begin{exercise}\label{ex:chi-pde-derivation}
  Derive equation~\eqref{eq:chi-pde} from~\eqref{eq:lindblad-cavity-vac}
  using the displacement-operator identities~\eqref{eq:D-deriv-identities}.
  Verify that the Gaussian ansatz yields~\eqref{eq:chi-solution}.
\end{exercise}
